\documentclass[aspectratio=169,11pt]{beamer}

% Packages
\usepackage[utf8]{inputenc}
\usepackage[T1]{fontenc}
\usepackage{lmodern}
\usepackage{microtype}
\usepackage{amsmath,amssymb}
\usepackage{booktabs}
\usepackage{graphicx}
\usepackage{tikz}
\usepackage{pgfplots}
\usepackage{xcolor}
\usepackage{ragged2e}
\usepackage{colortbl}
\usepackage{array}
\usepackage{hyperref}
\usepackage{adjustbox}

\usetikzlibrary{shapes.geometric, arrows.meta, positioning, calc, backgrounds, decorations.pathreplacing, shadows, fadings}
\pgfplotsset{compat=1.18}

% --- Color Palette: Warm Professional ---
\definecolor{DeepNavy}{HTML}{2E4057}
\definecolor{Teal}{HTML}{048A81}
\definecolor{WarmOrange}{HTML}{E85D04}
\definecolor{SoftPurple}{HTML}{9D4EDD}
\definecolor{WarmGray}{HTML}{6C757D}
\definecolor{LightGray}{HTML}{E9ECEF}
\definecolor{Cream}{HTML}{FBF8F1}
\definecolor{DeepRed}{HTML}{D62828}
\definecolor{Gold}{HTML}{D4A03A}
\definecolor{SoftWhite}{HTML}{FAFAFA}

% Beamer colors
\setbeamercolor{normal text}{fg=DeepNavy, bg=SoftWhite}
\setbeamercolor{structure}{fg=DeepNavy}
\setbeamercolor{alerted text}{fg=DeepRed}
\setbeamercolor{example text}{fg=Teal}
\setbeamercolor{frametitle}{fg=DeepNavy, bg=SoftWhite}
\setbeamercolor{title}{fg=DeepNavy}
\setbeamercolor{subtitle}{fg=WarmGray}
\setbeamercolor{block title}{bg=DeepNavy, fg=SoftWhite}
\setbeamercolor{block body}{bg=LightGray, fg=DeepNavy}
\setbeamercolor{itemize item}{fg=WarmOrange}
\setbeamercolor{itemize subitem}{fg=Gold}

% Fonts
\usefonttheme{professionalfonts}
\setbeamerfont{title}{size=\huge, series=\bfseries}
\setbeamerfont{frametitle}{size=\Large, series=\bfseries}

% Strip chrome
\setbeamertemplate{navigation symbols}{}
\setbeamertemplate{headline}{}

% Minimal footline
\setbeamertemplate{footline}{%
    \hfill
    \begin{beamercolorbox}[wd=3cm,ht=2.5ex,dp=1ex,right,rightskip=0.6cm]{page number}
        \usebeamercolor[fg]{WarmGray}\scriptsize\insertframenumber
    \end{beamercolorbox}
    \vspace{0.3cm}
}

% Bullets
\setbeamertemplate{itemize item}{\footnotesize\raisebox{0.3ex}{\tikz\fill[WarmOrange] (0,0) circle (0.45ex);}}
\setbeamertemplate{itemize subitem}{\footnotesize\raisebox{0.3ex}{\tikz\fill[Gold] (0,0) circle (0.35ex);}}

% Custom commands
\newcommand{\transitionslide}[2]{%
    {\setbeamercolor{normal text}{fg=SoftWhite,bg=DeepNavy}
    \begin{frame}[plain]\vfill\begin{center}
        {\huge\bfseries\textcolor{Gold}{#1}}\\[0.6cm]
        {\large\textcolor{LightGray}{#2}}
    \end{center}\vfill\end{frame}}
}

\begin{document}

\transitionslide{Technical Supplement}{Estimating Factor Loadings with PCA}

% --- SLIDE 1 ---
\begin{frame}
    \frametitle{The Goal: Recovering unobserved factor loadings}
    
    \vspace{0.3cm}
    Recall the firm-specific labor demand equation:
    \begin{equation*}
        L_{it} = \phi^d w_t + \lambda_i \eta_t + u_{it}
    \end{equation*}
    
    \vspace{0.3cm}
    \begin{itemize}
        \item \textbf{The Challenge:} We need weights $\Gamma$ orthogonal to factor loadings ($\sum \Gamma_i \lambda_i = 0$) to isolate idiosyncratic shocks. However, $\lambda_i$ (macro-sensitivity) is \textbf{unobserved}.
        \item \textbf{The Solution:} If idiosyncratic shocks $u_{it}$ are small, the dominant variance in labor demand $L_{it}$ is driven by the macro shock $\eta_t$.
        \item \textbf{The Method:} We can use \textbf{Principal Component Analysis (PCA)} on the panel of labor outcomes to extract the largest source of variance.
    \end{itemize}
\end{frame}

% --- SLIDE 2 ---
\begin{frame}
    \frametitle{Step 1: The Residual Matrix and Covariance}
    
    \vspace{0.5cm}
    Assume we partial out the wage effect $\phi^d w_t$. We are left with the residual labor variation driven by shocks:
    \begin{equation*}
        Y_{it} = \lambda_i \eta_t + u_{it}
    \end{equation*}
    
    \vspace{0.5cm}
    Let $Y$ be the $T \times N$ matrix of these demeaned observations. The sample \textbf{covariance matrix} across firms is:
    \begin{equation*}
        \Sigma = \frac{1}{T} Y'Y
    \end{equation*}
    
    \vspace{0.2cm}
    Because firms are primarily correlated through the common macro shock $\eta_t$, the matrix $\Sigma$ contains the structural signal of $\lambda$.
\end{frame}

% --- SLIDE 3 ---
\begin{frame}
    \frametitle{Step 2: Formulating the Optimization Problem}
    
    \vspace{0.5cm}
    {\Large PCA seeks a vector of weights $v \in \mathbb{R}^N$ that \textbf{maximizes variance}.}
    
    \vspace{0.8cm}
    The variance of the data projected onto $v$ is proportional to $v' \Sigma v$. To prevent infinite variance, we constrain the length of $v$ to 1.
    
    \vspace{0.5cm}
    \textbf{The Constrained Optimization Problem:}
    \begin{align*}
        \max_{v} \quad & v' \Sigma v \\
        \text{subject to} \quad & v'v = 1
    \end{align*}
\end{frame}

% --- SLIDE 4 ---
\begin{frame}[shrink=10]
    \frametitle{Step 3: Solving via Lagrange Multipliers}
    
    \vspace{0.5cm}
    We construct the Lagrangian $\mathcal{L}$ with a multiplier $\alpha$ for the constraint $v'v - 1 = 0$:
    \begin{equation*}
        \mathcal{L}(v, \alpha) = v' \Sigma v - \alpha (v'v - 1)
    \end{equation*}
    
    \vspace{0.5cm}
    Taking the derivative with respect to the vector $v$ and setting it to zero:
    \begin{align*}
        \frac{\partial \mathcal{L}}{\partial v} = 2 \Sigma v - 2 \alpha v &= 0 \\
        \Sigma v &= \alpha v
    \end{align*}
    
    \vspace{0.5cm}
    This is the fundamental \textbf{Eigenvalue Equation}. The solution $v$ is an eigenvector of the covariance matrix $\Sigma$, and $\alpha$ is its eigenvalue.
\end{frame}

% --- SLIDE 5 ---
\begin{frame}
    \frametitle{Step 4: Eigenvalues measure captured variance}
    
    \vspace{0.5cm}
    What does the eigenvalue $\alpha$ represent economically? Pre-multiply the eigenvalue equation by $v'$:
    
    \begin{align*}
        \Sigma v &= \alpha v \\
        v' \Sigma v &= v' (\alpha v) \\
        v' \Sigma v &= \alpha (v'v) \\
        v' \Sigma v &= \alpha \quad \text{(since } v'v = 1\text{)}
    \end{align*}
    
    \vspace{0.5cm}
    The eigenvalue $\alpha$ is \textbf{exactly equal} to the variance captured by the eigenvector $v$. Since we want to maximize variance, we select the eigenvector $v_1$ associated with the \textbf{largest} eigenvalue $\alpha_1$.
\end{frame}

% --- SLIDE 6 ---
\begin{frame}
    \frametitle{Step 5: Recovering the GIV Weights}
    
    \vspace{0.5cm}
    The principal eigenvector $v_1$ captures the dominant variance, which in our model is the macro shock $\eta_t$. Therefore, $v_1$ is our empirical estimate of the factor loadings:
    \begin{equation*}
        \hat{\lambda} = v_1
    \end{equation*}
    
    \vspace{0.5cm}
    With $\hat{\lambda}$ recovered, we can now construct the optimal Granular Instrument:
    \begin{enumerate}
        \item Compute the orthogonal projection matrix: $\hat{Q} = I - \hat{\lambda}(\hat{\lambda}'\hat{\lambda})^{-1}\hat{\lambda}'$
        \item Project the size weights: $\hat{\Gamma}^* = S'\hat{Q}$
        \item Construct the GIV: $z_t = \sum \hat{\Gamma}^*_i L_{it}$
    \end{enumerate}
\end{frame}

\end{document}
